\section{Data-bus arbiter (\texttt{dmem\_arb})}
\textit{Source: \texttt{design/interconnect/src/dmem\_arb.sv}.}

Three masters want the one data bus: the page-table walker (PTW), the atomic unit
(AMO), and the normal load/store unit (\sig{core2avl}, called \sig{c2a} here).
\sig{dmem\_arb} decides who gets the bus and routes the reply back to the right
master.

\subsection{Who wins}
The priority is fixed: \textbf{PTW $>$ AMO $>$ c2a}. The grants are
combinational:
\begin{itemize}[leftmargin=1.4em,itemsep=1pt]
\item \sig{ptw\_grant\_o}: PTW wins whenever \sig{ptw\_active\_i} is set;
\item \sig{amo\_grant\_o}: AMO wins when \sig{amo\_active\_i} is set, PTW is idle,
  and \sig{suppress\_data} is low;
\item \sig{c2a\_grant\_o}: c2a wins when neither PTW nor AMO is active and
  \sig{suppress\_data} is low.
\end{itemize}
\sig{suppress\_data} is the OR of \sig{mmu\_d\_access\_fault\_i},
\sig{mmu\_d\_fault\_i}, and \sig{d\_xlat\_pending\_i}. It blocks AMO and c2a (but
not PTW) while the MMU is
translating or a data fault is firing, so no access commits with a half-translated
or faulting address. The chosen master drives the shared slave port \sig{bus}
(\sig{bus.addr}, \sig{bus.we}, \sig{bus.be}, \sig{bus.wdata}, \sig{bus.req}). AMO
and c2a share the MMU-translated address \sig{data\_paddr\_i}; PTW uses its own
\sig{ptw\_addr\_i}.

\subsection{Routing the reply}
A read reply comes back later, so the arbiter must remember who asked.
\sig{pending\_rd\_master\_q} (one of \sig{M\_NONE}/\sig{M\_PTW}/\sig{M\_AMO}/\sig{M\_C2A})
records the owner of the outstanding read. It is set when a read is accepted
(\sig{read\_accepted = bus.req \& \textasciitilde bus.we \& bus.ready}) and cleared when the
reply arrives (\sig{bus.rvalid}). Then each master only sees its own reply:
\sig{c2a\_rvalid\_o = (pending\_rd\_master\_q == M\_C2A) \& bus.rvalid}, and the
same for PTW and AMO. \sig{*\_ready\_o = *\_grant\_o \& bus.ready}.

\latencynote{This ownership flop is what makes the bus safe under variable
latency. With single-cycle memory the reply came back the next cycle and could
not be confused with anyone else's. With a slow slave several requests and
replies overlap in time, so \sig{pending\_rd\_master\_q} is needed to send each
reply to the master that actually asked. It assumes one outstanding read per
master, which holds because each slave is blocking (it drops \sig{ready} during a
miss). The core pairs this with \sig{c2a\_load\_pending\_q}, which tracks the
same outstanding load on the CPU side.}
